% Detalle Técnico de la Implementación - Sección 5.3.1
% Para insertar en el Capítulo 5 de tu tesis

\subsection{Detalle Técnico de la Implementación}

Esta sección profundiza en los aspectos técnicos críticos del desarrollo de SGM-Payroll, explicando las decisiones de diseño, algoritmos implementados y optimizaciones aplicadas para garantizar la eficiencia y confiabilidad del sistema.

\subsubsection{Estructura y Diseño de la Base de Datos}

El modelo de datos de SGM-Payroll fue diseñado para optimizar las consultas de conciliación y garantizar la integridad referencial entre las diferentes fuentes de información.

\paragraph{Tablas Principales y Relaciones}

El sistema se estructura alrededor de cinco entidades centrales:

\begin{itemize}
    \item \textbf{CierreNomina}: Entidad principal que agrupa todos los archivos y resultados de un período específico
    \item \textbf{EmpleadoLibroRemuneraciones}: Datos base de empleados extraídos de Talana
    \item \textbf{MovimientosTalana}: Ausentismos, ingresos y finiquitos reportados por el sistema
    \item \textbf{ArchivoAnalista}: Datos ingresados manualmente por el analista
    \item \textbf{DiscrepanciasDetectadas}: Inconsistencias encontradas entre fuentes
\end{itemize}

\paragraph{Estrategia de Indexación}

Para optimizar el rendimiento de las consultas de conciliación, se implementaron índices estratégicos:

\begin{lstlisting}[language=SQL, caption=Índices Optimizados para Conciliación]
-- Índice compuesto para búsquedas por cliente y período
CREATE INDEX idx_cierre_cliente_periodo 
ON cierre_nomina (cliente_id, periodo, estado);

-- Índice único para empleados por RUT y cierre
CREATE UNIQUE INDEX idx_empleado_rut_cierre 
ON empleado_libro_remuneraciones (rut, cierre_nomina_id);

-- Índice para acelerar joins en comparaciones
CREATE INDEX idx_movimientos_empleado_fecha 
ON movimientos_talana (empleado_rut, fecha_movimiento, tipo_movimiento);

-- Índice para búsquedas de discrepancias por tipo
CREATE INDEX idx_discrepancias_tipo_estado 
ON discrepancias_detectadas (tipo_discrepancia, estado_resolucion);
\end{lstlisting}

\paragraph{Justificación de la Indexación}

La estrategia de indexación responde a los patrones de consulta más frecuentes del sistema:

\begin{itemize}
    \item \textbf{Filtrado por Cliente}: Cada consulta se filtra por cliente asignado (política de seguridad)
    \item \textbf{Búsquedas por RUT}: Operación más crítica para conciliación entre fuentes
    \item \textbf{Comparaciones Temporales}: Joins frecuentes por fecha para detectar inconsistencias
    \item \textbf{Agrupación por Tipo}: Clasificación de discrepancias para reportes ejecutivos
\end{itemize}

\subsubsection{Algoritmos de Comparación y Conciliación}

El sistema implementa tres tipos de comparaciones automatizadas para detectar discrepancias entre fuentes de datos.

\paragraph{Comparación de Movimientos de Personal}

Algoritmo para detectar inconsistencias entre movimientos reportados por Talana y el analista:

\begin{lstlisting}[language=SQL, caption=Detección de Discrepancias en Movimientos]
-- Empleados reportados como ingreso en Talana pero no en archivos analista
SELECT 
    t.empleado_rut,
    t.fecha_movimiento,
    'Ingreso en Talana no reportado por analista' as discrepancia
FROM movimientos_talana t
LEFT JOIN archivo_analista a 
    ON t.empleado_rut = a.empleado_rut 
    AND t.fecha_movimiento = a.fecha_movimiento
    AND t.tipo_movimiento = a.tipo_movimiento
WHERE 
    t.tipo_movimiento = 'INGRESO'
    AND a.empleado_rut IS NULL
    AND t.cierre_nomina_id = %s;

-- Diferencias en fechas de finiquito
SELECT 
    t.empleado_rut,
    t.fecha_movimiento as fecha_talana,
    a.fecha_movimiento as fecha_analista,
    ABS(EXTRACT(days FROM t.fecha_movimiento - a.fecha_movimiento)) as dias_diferencia
FROM movimientos_talana t
INNER JOIN archivo_analista a 
    ON t.empleado_rut = a.empleado_rut
WHERE 
    t.tipo_movimiento = 'FINIQUITO'
    AND t.fecha_movimiento != a.fecha_movimiento
    AND t.cierre_nomina_id = %s;
\end{lstlisting}

\paragraph{Comparación de Conceptos-Valor}

Para detectar diferencias en montos entre el Libro de Remuneraciones y las Novedades del analista:

\begin{lstlisting}[language=Python, caption=Algoritmo de Comparación de Conceptos]
def comparar_conceptos_empleado(empleado_rut, cierre_id):
    """
    Compara conceptos-valor entre fuentes para un empleado específico
    Implementa tolerancia de $100 para diferencias menores
    """
    tolerancia_monto = 100  # Pesos chilenos
    
    # Obtener conceptos del libro (Talana)
    conceptos_talana = get_conceptos_libro(empleado_rut, cierre_id)
    
    # Obtener novedades del analista
    novedades_analista = get_novedades_analista(empleado_rut, cierre_id)
    
    discrepancias = []
    
    for concepto, monto_talana in conceptos_talana.items():
        monto_analista = novedades_analista.get(concepto, 0)
        diferencia = abs(monto_talana - monto_analista)
        
        if diferencia > tolerancia_monto:
            discrepancias.append({
                'empleado_rut': empleado_rut,
                'concepto': concepto,
                'monto_talana': monto_talana,
                'monto_analista': monto_analista,
                'diferencia': diferencia,
                'tipo_discrepancia': 'DIFERENCIA_MONTO',
                'criticidad': 'ALTA' if diferencia > 10000 else 'MEDIA'
            })
    
    return discrepancias
\end{lstlisting}

\paragraph{Validación de Completitud de Datos}

Sistema para detectar empleados o conceptos faltantes entre fuentes:

\begin{lstlisting}[language=SQL, caption=Detección de Datos Faltantes]
-- Empleados en libro pero sin novedades del analista
WITH empleados_libro AS (
    SELECT DISTINCT empleado_rut 
    FROM empleado_libro_remuneraciones 
    WHERE cierre_nomina_id = %s
),
empleados_novedades AS (
    SELECT DISTINCT empleado_rut 
    FROM archivo_analista 
    WHERE cierre_nomina_id = %s AND tipo_archivo = 'NOVEDADES'
)
SELECT 
    el.empleado_rut,
    'Empleado sin novedades de analista' as discrepancia,
    'COMPLETITUD' as tipo_discrepancia
FROM empleados_libro el
LEFT JOIN empleados_novedades en ON el.empleado_rut = en.empleado_rut
WHERE en.empleado_rut IS NULL;
\end{lstlisting}

\subsubsection{Sistema de Mapeo de Conceptos Heterogéneos}

Una de las complejidades críticas del sistema es la normalización de conceptos entre diferentes fuentes que utilizan terminologías distintas.

\paragraph{Estrategia de Mapeo Automático}

El sistema implementa un mapeo jerárquico que combina reglas automáticas con intervención manual:

\begin{lstlisting}[language=Python, caption=Sistema de Mapeo Automático]
class MapeadorConceptos:
    def __init__(self, cliente_id):
        self.cliente_id = cliente_id
        self.mapeos_automaticos = self._cargar_mapeos_cliente()
        self.patrones_comunes = {
            r'.*sueldo.*base.*': 'SUELDO_BASE',
            r'.*afp.*': 'DESCUENTO_AFP',
            r'.*isapre.*': 'DESCUENTO_ISAPRE',
            r'.*horas.*extra.*50.*': 'HORAS_EXTRA_50',
            r'.*gratific.*': 'GRATIFICACION'
        }
    
    def mapear_concepto(self, concepto_original):
        """
        Mapeo jerárquico: explícito -> patrón -> manual pendiente
        """
        # 1. Mapeo explícito por cliente
        if concepto_original in self.mapeos_automaticos:
            return self.mapeos_automaticos[concepto_original]
        
        # 2. Mapeo por patrones regex
        for patron, concepto_estandar in self.patrones_comunes.items():
            if re.match(patron, concepto_original.lower()):
                # Guardar mapeo encontrado para futuros usos
                self._guardar_mapeo_automatico(concepto_original, concepto_estandar)
                return concepto_estandar
        
        # 3. Concepto requiere mapeo manual
        self._marcar_mapeo_pendiente(concepto_original)
        return None
    
    def obtener_mapeos_pendientes(self):
        """Retorna conceptos que requieren intervención manual"""
        return ConceptosPendientesMapeo.objects.filter(
            cliente_id=self.cliente_id, 
            estado='PENDIENTE'
        ).values_list('concepto_original', flat=True)
\end{lstlisting}

\paragraph{Persistencia de Mapeos}

Los mapeos se almacenan en una estructura optimizada para consultas frecuentes:

\begin{lstlisting}[language=SQL, caption=Tabla de Mapeos con Índices]
CREATE TABLE mapeo_concepto (
    id SERIAL PRIMARY KEY,
    cliente_id INTEGER NOT NULL REFERENCES cliente(id),
    concepto_original VARCHAR(255) NOT NULL,
    concepto_estandar VARCHAR(100) NOT NULL,
    tipo_mapeo VARCHAR(20) NOT NULL, -- 'AUTOMATICO', 'MANUAL', 'PATRON'
    activo BOOLEAN DEFAULT TRUE,
    fecha_creacion TIMESTAMP DEFAULT NOW(),
    creado_por INTEGER REFERENCES usuario(id)
);

-- Índice único para evitar duplicados y acelerar búsquedas
CREATE UNIQUE INDEX idx_mapeo_cliente_concepto 
ON mapeo_concepto (cliente_id, concepto_original) WHERE activo = TRUE;

-- Índice para consultas por tipo de mapeo
CREATE INDEX idx_mapeo_tipo 
ON mapeo_concepto (tipo_mapeo, activo);
\end{lstlisting}

\subsubsection{Optimización de Rendimiento}

\paragraph{Procesamiento Asíncrono con Celery}

Para manejar archivos de gran tamaño (133 empleados × 55 conceptos), el sistema implementa un pipeline asíncrono:

\begin{lstlisting}[language=Python, caption=Pipeline de Procesamiento Asíncrono]
@shared_task(bind=True)
def procesar_archivo_completo(self, cierre_id, archivo_path):
    """
    Pipeline completo de procesamiento por fases
    Permite recovery en caso de fallos
    """
    try:
        # Fase 1: Análisis de headers (rápido)
        resultado_headers = analizar_headers.delay(cierre_id, archivo_path)
        
        # Fase 2: Validación de datos (medio)
        resultado_validacion = validar_datos.delay(cierre_id)
        
        # Fase 3: Procesamiento de empleados (lento)
        resultado_empleados = procesar_empleados.delay(cierre_id)
        
        # Fase 4: Generación de discrepancias (crítico)
        resultado_discrepancias = detectar_discrepancias.delay(cierre_id)
        
        # Chain permite rollback granular por fase
        chain = (resultado_headers | resultado_validacion | 
                resultado_empleados | resultado_discrepancias)
        
        return chain.apply_async()
        
    except Exception as e:
        # Log estructurado para debugging
        logger.error(f"Error en procesamiento cierre {cierre_id}: {str(e)}")
        self.retry(countdown=60, max_retries=3)
\end{lstlisting}

\paragraph{Cache Estratégico con Redis}

Sistema de cache optimizado por patrones de acceso:

\begin{lstlisting}[language=Python, caption=Implementación de Cache Estratégico]
class CacheEstrategico:
    def __init__(self):
        self.redis_client = redis.Redis(db=2)  # DB dedicado para nómina
        self.ttl_por_tipo = {
            'mapeos': 86400,      # 24 horas - críticos
            'empleados': 7200,    # 2 horas - datos temporales
            'discrepancias': 3600, # 1 hora - resultados volátiles
            'kpis': 21600         # 6 horas - reportes ejecutivos
        }
    
    def get_key_estructurada(self, cliente_id, cierre_id, componente, tipo):
        """Genera keys organizadas por patrón estándar"""
        return f"sgm:nomina:{cliente_id}:{cierre_id}:{componente}:{tipo}"
    
    def cache_mapeos_empleado(self, cliente_id, cierre_id, empleado_rut, conceptos):
        """Cache selectivo solo para empleados en procesamiento activo"""
        key = self.get_key_estructurada(cliente_id, cierre_id, 
                                       f"empleado:{empleado_rut}", "conceptos")
        
        self.redis_client.setex(key, self.ttl_por_tipo['empleados'], 
                               json.dumps(conceptos))
    
    def limpiar_por_fase(self, cliente_id, cierre_id, fase):
        """Cleanup automático al completar fases"""
        pattern = f"sgm:nomina:{cliente_id}:{cierre_id}:{fase}:*"
        keys = self.redis_client.keys(pattern)
        
        if keys:
            self.redis_client.delete(*keys)
            logger.info(f"Limpieza automática: {len(keys)} keys eliminadas de {fase}")
\end{lstlisting}

\subsubsection{Validaciones y Control de Integridad}

\paragraph{Validación de RUT Chileno}

Implementación del algoritmo oficial de validación:

\begin{lstlisting}[language=Python, caption=Validador de RUT Optimizado]
def validar_rut_chileno(rut_completo):
    """
    Valida RUT usando algoritmo módulo 11
    Optimizado para procesamiento masivo
    """
    # Limpiar formato: '12.345.678-9' -> '123456789'
    rut_limpio = re.sub(r'[^0-9kK]', '', str(rut_completo))
    
    if len(rut_limpio) < 2:
        return False
    
    # Separar número y dígito verificador
    rut_numero = rut_limpio[:-1]
    dv_esperado = rut_limpio[-1].upper()
    
    # Algoritmo módulo 11
    suma = 0
    multiplicador = 2
    
    for digito in reversed(rut_numero):
        suma += int(digito) * multiplicador
        multiplicador = 3 if multiplicador == 7 else multiplicador + 1
    
    resto = suma % 11
    dv_calculado = 'K' if resto == 1 else ('0' if resto == 0 else str(11 - resto))
    
    return dv_esperado == dv_calculado

# Validación masiva optimizada con set comprehension
def validar_ruts_masivo(lista_ruts):
    """Validación optimizada para archivos con cientos de empleados"""
    ruts_invalidos = {rut for rut in lista_ruts if not validar_rut_chileno(rut)}
    return len(ruts_invalidos) == 0, ruts_invalidos
\end{lstlisting}

\paragraph{Control de Integridad de Archivos}

Sistema de validaciones en cascada antes del procesamiento:

\begin{lstlisting}[language=Python, caption=Validaciones de Integridad]
class ValidadorArchivos:
    def validar_libro_remuneraciones(self, archivo_path):
        """Validaciones específicas para libro de Talana"""
        errores = []
        
        try:
            df = pd.read_excel(archivo_path)
            
            # Validación 1: Estructura mínima
            if df.empty:
                errores.append("Archivo vacío")
                return False, errores
            
            # Validación 2: Columnas obligatorias
            columnas_requeridas = ['rut', 'nombre', 'apellido']
            columnas_faltantes = [col for col in columnas_requeridas 
                                if not any(col.lower() in c.lower() for c in df.columns)]
            
            if columnas_faltantes:
                errores.append(f"Columnas faltantes: {columnas_faltantes}")
            
            # Validación 3: RUTs duplicados
            rut_col = self._encontrar_columna_rut(df.columns)
            if rut_col:
                ruts_duplicados = df[df[rut_col].duplicated()][rut_col].tolist()
                if ruts_duplicados:
                    errores.append(f"RUTs duplicados encontrados: {len(ruts_duplicados)}")
            
            # Validación 4: Valores numéricos en conceptos
            conceptos_numericos = self._detectar_conceptos_numericos(df)
            for concepto in conceptos_numericos:
                valores_invalidos = df[~pd.to_numeric(df[concepto], errors='coerce').notna()]
                if not valores_invalidos.empty:
                    errores.append(f"Valores no numéricos en {concepto}: {len(valores_invalidos)} filas")
            
            return len(errores) == 0, errores
            
        except Exception as e:
            return False, [f"Error procesando archivo: {str(e)}"]
\end{lstlisting}

Esta implementación técnica detallada demuestra la robustez y sofisticación del sistema SGM-Payroll, evidenciando decisiones de diseño fundamentadas en los requerimientos específicos del procesamiento de nóminas empresariales.
